\documentclass[10pt,a4paper]{amsart}
\usepackage[margin=19mm]{geometry}
\usepackage{amsmath,amssymb,mathtools}
\usepackage[scr=rsfs,cal=euler]{mathalfa}
\usepackage{xfrac}
\usepackage[unicode,hidelinks,bookmarks=false]{hyperref}
\newcommand{\ie}{i.e.\ }
\newcommand{\cf}{cf.\ }
\newcommand{\resp}{respectively\ }
\newcommand{\Subsection}{\subsection}
\providecommand{\nobreakdash}{\nobreak\hbox{-}}
\setlength{\emergencystretch}{2em}
\multlinegap0pt
\usepackage{bm}
\usepackage{bbm}
\usepackage{dsfont}
\usepackage{xspace}
\usepackage{cancel}
\usepackage{mathtools}
\usepackage{amsthm}
\makeatletter
\@ifundefined{theorem}{\newtheorem{theorem}{Theorem}}{}
\@ifundefined{proposition}{\newtheorem{proposition}[theorem]{Proposition}}{}
\makeatother
\def\balpha   {\boldsymbol{\alpha}}
\def\cM { \mathcal{M}}
\def\cR { \mathcal{R}}
\def\grad            {\textrm{grad}}     % grad   
\def\v {\bm{v}}
\def\x {\bm{x}} 
\def\zeros           {\bm{0}}            % All zero vector or matrix
\title[Chordal-NMF with Riemannian Multiplicative Update]{Chordal-NMF with Riemannian Multiplicative Update}
\author{Flavia Esposito, Andersen Ang}
\date{Source: arXiv:2405.12823v4 (CC BY 4.0)}
\begin{document}
\maketitle

\noindent\emph{Source and licence.} Chordal-NMF with Riemannian Multiplicative Update. Version arXiv:2405.12823v4 (CC BY 4.0), distributed under the Creative Commons Attribution 4.0 International licence, as recorded in the arXiv licence metadata for this exact version and shown on the abstract page https://arxiv.org/abs/2405.12823v4. Full rights notice: licenses/arxiv-2405.12823-CC-BY-4.0.txt.\par\smallskip

\noindent\textit{Editorial scope.} Counted unit: the Proposition whose source label is \texttt{prop:RMU\_update} in the article (source0.tex lines 780--784, source label \texttt{prop:RMU\_update}), delivered together with the article's own complete original proof of it (source0.tex lines 785--815, one \texttt{proof} environment of 31 lines). No mathematics is written, repaired or re-derived: statement and proof are the article's own text line for line. Required context retained uncounted: metric\_retraction (lines 699--703). Every definition, display and label of those lines is retained unchanged. Omitted from this package and not redistributed: 1--698, 704--779, 816--1915. Nothing inside a retained excerpt is omitted or altered: every retained body is delivered exactly as the article's lines are written, so no omission receipt of any kind accompanies this package. Citations: the retained excerpts carry no citation command at all, so no bibliography entry of the article is transcribed and no reference is redistributed. Reference adaptation: none was needed and none was made; every reference inside the delivered unit resolves inside it, so no unresolved reference or citation is produced. Printed slips kept literal: no printed word, symbol, hypothesis or number of the counted statement or of its proof was altered. Layout and class: the article's own class is \texttt{sn-jnl} with packages including mathtools, natbib, graphicx, multirow, amsmath, amssymb, amsfonts, amsthm, mathrsfs, appendix, xcolor, textcomp, manyfoot, booktabs; the wrapper is the lane's pinned \texttt{amsart} editorial wrapper at 10pt on A4, which restates the theorem-like environments the retained text needs and adds the article's own macro definitions that the retained text needs (\texttt{\textbackslash balpha}, \texttt{\textbackslash cM}, \texttt{\textbackslash cR}, \texttt{\textbackslash grad}, \texttt{\textbackslash v}, \texttt{\textbackslash x}, \texttt{\textbackslash zeros}), with extra wrapper packages (bm, bbm, dsfont, xspace, cancel, mathtools). The delivered numbering is the unit's own. Assets: no retained span contains a figure environment, a graphics-inclusion command, a PDF-page inclusion, a table or a TikZ picture; no image file is copied into this package, the package declares no asset dependency and no image file is redistributed with it. Licence: version arXiv:2405.12823v4 of this article is distributed under the Creative Commons Attribution 4.0 International licence, as recorded in the arXiv licence metadata for this exact version and shown on the abstract page https://arxiv.org/abs/2405.12823v4. A full rights notice is delivered at licenses/arxiv-2405.12823-CC-BY-4.0.txt. Characters: every retained excerpt is pure ASCII, so the delivered engine, XeLaTeX with its default Latin Modern text font, reports no missing character.
\begin{equation}\label{metric_retraction}
\cR_{\x}(\v) \coloneqq 
\dfrac{\x+\v}{\|\x+\v\|_E}.
\tag{Metric retraction}
\end{equation}

\begin{proposition}[RMU]\label{prop:RMU_update}
Denote $\v_k = - \grad f(\x_k)$ the anti-parallel direction of the Riemannian gradient of a manifold $\cM$ at $\x_k$, and let $\cR_{\x_k}$ be the metric retraction onto $\cM$.
If a nonnegative $\x_k$ is updated by $\x_{k+1} = \cR_{\x_k}( \balpha \odot  \v_k)$
with an element-wise stepsize $\balpha=\x_k \oslash \grad^+ f(\x_k) \in E$, then $\x_{k+1} \geq \zeros$ and is on $\cM$.
\end{proposition}

    \begin{proof}
The term $\balpha \odot  (-\grad f(\x_k))$ with $\balpha=\x_k \oslash \grad^+ f(\x_k)$ can be computed as 
\[
\begin{array}{lll}
-\balpha \odot  \grad f(\x_k) 
&=
\big( \x_k \oslash \grad^+ f(\x_k) \big)  \odot  \big( \grad^- f(\x_k) - \grad^+ f(\x_k) \big)
\\
&=
\x_k \odot \grad^- f(\x_k) \oslash \grad^+ f(\x_k) - \x_k .
\end{array}
\]
The element-wise operations as linear transformations  in general do not preserve the tangent condition of a manifold into a point.
However here we can still apply~\eqref{metric_retraction} on $-\balpha \odot  \grad f(\x_k)$, the $\x_k$ terms got canceled and gives
\[
\begin{array}{ll}
\x_{k+1} 
=
\cR_{\x_k}\big( -\balpha \odot  \grad f(\x_k)\big) 
&=
\dfrac{\x_k - \balpha \odot  \grad f(\x_k)}{\| \x_k - \balpha \odot  \grad f(\x_k) \|_E}
\\
&= 
\dfrac{ \x_k \odot \grad^- f(\x_k) \oslash \grad^+ f(\x_k)}
{
\|  \x_k \odot \grad^- f(\x_k) \oslash \grad^+ f(\x_k) \|_E
}.    
\end{array}
\]
The ratio between $\grad^+ f$ and $\grad^- f$ is nonnegative as its parts are nonnegative by definition, the denominator is nonnegative, therefore $\x_{k+1}$ is nonnegative if $\x_k$ is nonnegative.
\end{proof}

\end{document}
